% =====================================================================
% Results
% =====================================================================
\section{Results}
\label{sec:results}

\subsection{Baseline scenarios}
\label{sec:results-baseline}

\Cref{fig:mobile-traj} shows the four interpolated trajectories for
the 12-waypoint mobile-robot baseline. Linear interpolation is
identical to the waypoint polygon by construction. The cubic and
quintic splines pass exactly through every waypoint and round the
two 90$^{\circ}$ corners with visibly different tightness. The
B-spline approximation deviates noticeably from the interior
waypoints, cutting both 90$^{\circ}$ corners and producing a
shorter overall path.

\begin{figure}[!t]
\centering
\includegraphics[width=\linewidth]{mobile_robot_trajectory.png}
\caption{Mobile-robot trajectory: comparison of the four
interpolation methods on the 12-waypoint baseline. White circles are
the waypoints; numbered for reference.}
\label{fig:mobile-traj}
\end{figure}

The kinematic profile in \cref{fig:mobile-kinematics} reveals the
qualitative differences between the methods that are not visible from
the trajectory plot alone. Linear interpolation has zero analytic
acceleration and zero analytic jerk because each segment is a
constant-velocity straight line. Cubic spline jerk exhibits the
characteristic step pattern of a piecewise-cubic interpolant whose
third derivative is constant on each segment. Quintic spline jerk is
continuous but has visibly larger amplitude near $t=0$ and $t=T$,
where natural boundary conditions force a rapid transition from rest
to nominal speed. The B-spline produces the smoothest derivative
profile by every metric.

\begin{figure*}[!t]
\centering
\includegraphics[width=0.78\linewidth]{mobile_robot_kinematics.png}
\caption{Mobile-robot baseline: speed, acceleration magnitude, and
jerk magnitude profiles for each method. Note the piecewise-constant
jerk pattern of the cubic spline and the endpoint jerk spikes of the
natural-BC quintic spline.}
\label{fig:mobile-kinematics}
\end{figure*}

For the robot-arm scenario, \cref{fig:arm-cartesian} shows the
end-effector trajectories produced by mapping each method's
joint-space output through forward kinematics. Even though every
interpolating method visits the same eight joint waypoints, the
Cartesian paths between them differ because the kinematic mapping is
nonlinear. The natural-BC quintic produces visible loops near
waypoints 2 and 5--6 that are not present in the cubic.

\begin{figure}[!t]
\centering
\includegraphics[width=\linewidth]{robot_arm_cartesian.png}
\caption{Robot-arm Cartesian end-effector trajectories: result of
running each method on the joint-space waypoints and projecting
through forward kinematics.}
\label{fig:arm-cartesian}
\end{figure}

\begin{table*}[t]
\centering
\caption{Baseline metrics for all four methods on both scenarios.
Mobile-robot units are metres; arm units are radians for joint-space
quantities. \texttt{wp\_err} is the maximum waypoint deviation;
\texttt{disc} is the maximum velocity discontinuity at internal knots.}
\label{tab:baseline-metrics}
\begin{tabular}{llrrrrrrr}
\toprule
\textbf{Dataset} & \textbf{Method} & $\max\|v\|$ & $\max\|a\|$ & $\mathrm{rms}\|j\|$ & path & wp\_err & disc & $t$ (ms) \\
\midrule
Mobile robot & Linear   & 1.000 & 0.000 & 0.000 & 24.93 & $0$ & $1.41$ & 0.13 \\
             & Cubic    & 1.215 & 1.175 & 0.340 & 25.51 & $0$ & $9{\times}10^{-5}$ & 0.37 \\
             & Quintic  & 1.709 & 1.491 & 0.854 & 25.52 & $0$ & $8{\times}10^{-5}$ & 0.75 \\
             & B-spline & 1.509 & 0.411 & 0.126 & 22.92 & $0.68$ & $3{\times}10^{-5}$ & 0.65 \\
\midrule
Robot arm    & Linear   & 1.888 & 0.000 & 0.000 & 6.37 & $0$ & --- & 0.12 \\
             & Cubic    & 2.580 & 4.225 & 5.191 & 6.66 & $0$ & --- & 0.29 \\
             & Quintic  & 2.958 & 6.810 & 10.986 & 7.63 & $0$ & --- & 0.67 \\
             & B-spline & 2.832 & 3.280 & 1.566 & 4.51 & $0.40$ & --- & 0.44 \\
\bottomrule
\end{tabular}
\end{table*}

The quantitative metrics in \cref{tab:baseline-metrics} confirm
several observations from the figures. The interpolating methods all
hit every waypoint to floating-point precision; only the B-spline
deviates (0.68\,m on the mobile robot, 0.40\,rad on the robot arm).
Linear interpolation reports zero analytic acceleration but a $\sim 1$\,m/s
velocity-discontinuity magnitude on the mobile robot, in contrast to
the smooth methods whose discontinuity values sit at $10^{-5}$\,m/s
and below (the latter is dominated by $\varepsilon$-proportional
numerical noise from the finite-difference probe). Surprisingly, the
quintic spline produces \emph{higher} peak acceleration and RMS jerk
than the cubic on both datasets despite its higher continuity class;
this anomaly is investigated in Section~\ref{sec:results-bc-sens}.

\subsection{Density sweep}
\label{sec:results-density}

\Cref{fig:density-sweep} sweeps over waypoint count
$n \in \{5, 12, 20\}$ on the mixed warehouse-aisle geometry under
chord-length parameterisation. Counter to a naive expectation that
``more waypoints means smoother motion'', peak acceleration and RMS
jerk both increase monotonically with $n$ for every spline-style
method.

\begin{figure*}[!t]
\centering
\includegraphics[width=\linewidth]{sweep_density.png}
\caption{Density sweep: peak acceleration, RMS jerk, and path length
versus waypoint count on the mixed mobile-robot geometry with
chord-length spacing. Acceleration and jerk increase monotonically;
path length converges towards the polyline.}
\label{fig:density-sweep}
\end{figure*}

The mechanism is straightforward when viewed alongside the
trajectories themselves (\cref{fig:density-trajectories}): denser
waypoints under chord-length spacing reduce the segment width $h_i$
proportionally, forcing larger curvatures into the splines and
generating proportionally larger time derivatives. Path length
converges towards the polyline as $n$ grows, demonstrating that
geometric fidelity improves while smoothness deteriorates.

\begin{figure*}[!t]
\centering
\includegraphics[width=\linewidth]{sweep_density_trajectories.png}
\caption{Mobile-robot trajectories at varying waypoint density. At
$n=5$ all smooth methods make obvious geometric mistakes; at $n=20$
they trace the polyline closely.}
\label{fig:density-trajectories}
\end{figure*}

\subsection{Geometry sweep}
\label{sec:results-geometry}

\Cref{fig:geometry-sweep} compares the sharp 90$^{\circ}$ zigzag and
the smooth sinusoidal scenarios at $n=12$, with both using
chord-length spacing. The sharp geometry approximately doubles the
peak acceleration and triples the RMS jerk relative to the gradual
sinusoid for cubic and quintic splines. The B-spline absorbs the
geometric character much more gracefully: its sharp-vs.-gradual ratio
for peak acceleration is only 1.18 (compared to 1.78 for the cubic).

\begin{figure*}[!t]
\centering
\includegraphics[width=\linewidth]{sweep_geometry.png}
\caption{Geometry sweep at $n=12$: sharp 90$^{\circ}$ zigzag versus
smooth sinusoid. The third panel uses a symmetric-log scale and
cleanly separates the linear interpolant ($\sim 1$\,m/s real
velocity discontinuity) from the smooth methods ($\sim 10^{-5}$\,m/s).}
\label{fig:geometry-sweep}
\end{figure*}

The trajectory plots in \cref{fig:geometry-trajectories} show why.
On the sharp zigzag, the cubic and quintic interpolants exhibit
Runge-style overshoot at every internal corner, attempting to
reconcile pass-through-every-waypoint with smooth derivatives. The
B-spline does not overshoot but cuts each corner aggressively, with
correspondingly small derivative magnitudes. On the smooth sinusoid,
all three smooth methods agree closely.

\begin{figure*}[!t]
\centering
\includegraphics[width=\linewidth]{sweep_geometry_trajectories.png}
\caption{Mobile-robot trajectories on the sharp zigzag (left) and
gradual sinusoid (right) variants at $n=12$. Cubic and quintic
splines overshoot at sharp corners; the B-spline cuts them.}
\label{fig:geometry-trajectories}
\end{figure*}

\subsection{Spacing sweep}
\label{sec:results-spacing}

The third axis is the smallest in magnitude
(\cref{fig:spacing-sweep}). Switching from chord-length to uniform
time spacing on the same 12-waypoint mixed geometry leaves peak
acceleration and RMS jerk within roughly 30\% across all methods.
The most visible effect is on peak speed: the linear interpolant's
chord-length-parameterised speed envelope is identically 1\,m/s, but
under uniform spacing it must traverse the long diagonal segments in
the same one-second budget as the short ones, pushing peak speed to
1.32\,m/s. The smooth methods inherit some of this effect to a
lesser degree.

\begin{figure*}[!t]
\centering
\includegraphics[width=\linewidth]{sweep_spacing.png}
\caption{Spacing sweep: chord-length versus uniform time spacing at
$n=12$ on the mixed geometry. Peak speed is the most affected
metric; peak acceleration and RMS jerk are within 30\%.}
\label{fig:spacing-sweep}
\end{figure*}

\subsection{Derivative validation}
\label{sec:results-validation}

\Cref{fig:validation} compares each method's analytic derivative
output against centered finite-difference estimates of the same
underlying dense position trajectory, using the schemes from
Section~\ref{sec:bg-diff} on the cubic spline run. The two curves
overlap to plotting precision in the velocity and acceleration
panels. The jerk panel shows the only meaningful disagreement: the
cubic spline has true step discontinuities in jerk at every internal
knot, and the centered five-point finite-difference scheme cannot
resolve those steps, instead smoothing each into a short ramp.

\begin{figure*}[!t]
\centering
\includegraphics[width=0.78\linewidth]{validation_analytic_vs_empirical.png}
\caption{Analytic derivative vs.\ centered finite-difference
empirical derivative for the cubic spline on the baseline mobile
robot. Velocity and acceleration overlap to plotting precision; the
jerk panel shows the FD scheme's expected smoothing of the
analytic step discontinuities at internal knots.}
\label{fig:validation}
\end{figure*}

The maximum interior errors are summarised in
\cref{tab:validation-errors}. The cubic-spline jerk error is
several orders of magnitude larger than the velocity or acceleration
errors because the latter are smooth functions of the cubic
polynomial pieces while the former is genuinely discontinuous. The
quintic and B-spline jerk errors are at the $10^{-3}$ to $10^{-1}$
level, consistent with their respective continuity classes.

\begin{table}[!t]
\centering
\caption{Maximum interior errors of analytic vs.\ finite-difference
derivatives, baseline mobile-robot dataset. Boundary stencils are
excluded from the comparison.}
\label{tab:validation-errors}
\begin{tabular}{lccc}
\toprule
\textbf{Method} & $\max |v_a - v_e|$ & $\max |a_a - a_e|$ & $\max |j_a - j_e|$ \\
\midrule
Cubic    & $2.6 \times 10^{-5}$ & $2.3 \times 10^{-3}$ & $5.1 \times 10^{-1}$ \\
Quintic  & $7.3 \times 10^{-8}$ & $4.3 \times 10^{-4}$ & $8.8 \times 10^{-4}$ \\
B-spline & $5.1 \times 10^{-6}$ & $6.6 \times 10^{-4}$ & $1.3 \times 10^{-1}$ \\
\bottomrule
\end{tabular}
\end{table}

\subsection{Quintic boundary-condition sensitivity}
\label{sec:results-bc-sens}

The boundary-condition sensitivity study
(Section~\ref{sec:bc-sensitivity-method}) identifies the natural BC
as the source of the quintic spline's anomalously high derivative
norms (\cref{fig:quintic-bc}). Replacing the natural BC
($\vec{v}_0 = \vec{a}_0 = \vec{v}_n = \vec{a}_n = \vec{0}$) with the
cubic-derived clamped BC reduces peak acceleration from
1.49\,m/s$^2$ to 1.00\,m/s$^2$ (a 33\% reduction) and RMS jerk from
0.85\,m/s$^3$ to 0.31\,m/s$^3$ (a 64\% reduction).

\begin{figure*}[!t]
\centering
\includegraphics[width=0.78\linewidth]{quintic_bc_sensitivity.png}
\caption{Quintic spline boundary-condition sensitivity, mobile-robot
baseline. Natural BC forces start/stop-at-rest in both velocity and
acceleration, generating large endpoint-region jerk. Clamped BC
with cubic-derived end-velocities and end-accelerations removes the
endpoint spikes entirely.}
\label{fig:quintic-bc}
\end{figure*}

With the cubic-derived clamped BC, the quintic is the smoothest
interpolating method on the baseline---outperforming the cubic on
every metric except peak velocity. The natural-BC quintic should
therefore be reported as a method-plus-BC choice in published
results, not as an inherent property of quintic splines.
